%====================================================%
% SEZIONE 1: Composizione di funzioni
%====================================================%

\begin{question}[ID=1.1, subtitle={Composizione di funzioni}, difficulty=facile, tags={funzioni, composizione}]
Date le funzioni $f(x)=2x-1$ e $g(x)=x+3$, calcola:
\begin{multicols}{2}
\begin{enumerate}[leftmargin=3.2cm]
    \item $(f\circ g)(2)$
    \item $(g\circ f)(2)$
    \item $(f\circ g)(x)$
    \item $(g\circ f)(x)$
    \item $(f\circ g)(-1)$
    \item Determina $x$ tale che $(f\circ g)(x)=7$
\end{enumerate}
\end{multicols}
\begin{solution}
\begin{enumerate}
    \item $9$
    \item $6$
    \item $(f\circ g)(x)=2x+5$
    \item $(g\circ f)(x)=2x+2$
    \item $3$
    \item $x=1$
\end{enumerate}
\end{solution}
\end{question}

%----------------------------------------------------%

\begin{question}[ID=1.2, subtitle={Composizione di funzioni polinomiali}, difficulty=facile, tags={funzioni, composizione}]
Date le funzioni $f(x)=x^2+1$ e $g(x)=2x-3$, calcola:
\begin{multicols}{2}
\begin{enumerate}[leftmargin=3.2cm]
    \item $(f\circ g)(x)$
    \item $(g\circ f)(x)$
    \item $(f\circ g)(2)$
    \item $(g\circ f)(-2)$
    \item Determina $x$ tale che $(f\circ g)(x)=10$
    \item Determina $x$ tale che $(g\circ f)(x)=17$
\end{enumerate}
\end{multicols}
\begin{solution}
\begin{enumerate}
    \item $(f\circ g)(x)=4x^2-12x+10$
    \item $(g\circ f)(x)=2x^2-1$
    \item $(f\circ g)(2)=2$
    \item $(g\circ f)(-2)=7$
    \item $x=0\vee x=3$
    \item $x=\pm3$
\end{enumerate}
\end{solution}
\end{question}

%----------------------------------------------------%

\begin{question}[ID=1.3, subtitle={Composizione di funzioni con frazioni}, difficulty=media, tags={funzioni, composizione}]
Date le funzioni $f(x)=\dfrac{x+1}{2}$ e $g(x)=x^2-1$, calcola:
\begin{multicols}{2}
\begin{enumerate}[leftmargin=3.2cm]
    \item $(f\circ g)(x)$
    \item $(g\circ f)(x)$
    \item $(f\circ g)(3)$
    \item $(g\circ f)(1)$
    \item Determina $x$ tale che $(f\circ g)(x)=8$
    \item Determina $x$ tale che $(g\circ f)(x)=6$
\end{enumerate}
\end{multicols}
\begin{solution}
\begin{enumerate}
    \item $(f\circ g)(x)=\dfrac{x^2}{2}$
    \item $(g\circ f)(x)=\dfrac{x^2+2x-3}{4}$
    \item $(f\circ g)(3)=\dfrac{9}{2}$
    \item $(g\circ f)(1)=0$
    \item $x=\pm4$
    \item $x=-1\pm\sqrt{27}$
\end{enumerate}
\end{solution}
\end{question}

%----------------------------------------------------%

\begin{question}[ID=1.4, subtitle={Composizione di funzioni con radicali}, difficulty=media, tags={funzioni, composizione}]
Date le funzioni $f(x)=\sqrt{2x+1}$ e $g(x)=x^2+1$, calcola:
\begin{multicols}{2}
\begin{enumerate}[leftmargin=3.2cm]
    \item $(f\circ g)(x)$
    \item $(g\circ f)(x)$
    \item $(f\circ g)(2)$
    \item $(g\circ f)(4)$
    \item Determina $x$ tale che $(g\circ f)(x)=8$
    \item Determina $x$ tale che $(f\circ g)(x)=\sqrt{11}$
\end{enumerate}
\end{multicols}
\begin{solution}
\begin{enumerate}
    \item $(f\circ g)(x)=\sqrt{2x^2+3}$
    \item $(g\circ f)(x)=2x+2$
    \item $(f\circ g)(2)=\sqrt{11}$
    \item $(g\circ f)(4)=10$
    \item $x=3$
    \item $x=\pm2$
\end{enumerate}
\end{solution}
\end{question}

%----------------------------------------------------%

\begin{question}[ID=1.5, subtitle={Composizione di tre funzioni}, difficulty=difficile, tags={funzioni, composizione}]
Date le funzioni $f(x)=2x-1$, $g(x)=x^2+1$ e $h(x)=3x-2$, calcola:
\begin{multicols}{2}
\begin{enumerate}[leftmargin=3.2cm]
    \item $(f\circ g\circ h)(x)$
    \item $(g\circ f\circ h)(x)$
    \item $(h\circ g\circ f)(x)$
    \item $(f\circ g)(h(1))$
    \item $(g\circ h)(1)$
    \item Determina $x$ tale che $(f\circ g\circ h)(x)=9$
\end{enumerate}
\end{multicols}
\begin{solution}
\begin{enumerate}
    \item $(f\circ g\circ h)(x)=18x^2-24x+9$
    \item $(g\circ f\circ h)(x)=36x^2-60x+26$
    \item $(h\circ g\circ f)(x)=12x^2-12x+4$
    \item $(f\circ g)(h(1))=19$
    \item $(g\circ h)(1)=2$
    \item $x=0\vee x=\tfrac43$
\end{enumerate}
\end{solution}
\end{question}

